\documentclass[12pt]{article}
\usepackage{graphicx}
\usepackage{amsmath}
\usepackage{amssymb}
\usepackage{url}
\usepackage{booktabs}
\usepackage[margin=1in]{geometry}
\usepackage{natbib}
\usepackage{setspace}
\usepackage{enumitem}

\title{Prohibition Next Door: Distance, Population Exposure, and Dispensary Revenue at U.S. Cannabis-Market Borders \\
{\large (Working title -- see Section~\ref{sec:titles} for alternative titles)}}
\author{[AUTHOR NAME(S)] \\ {}[AFFILIATION(S)]}
\date{[INSERT MONTH YEAR]}

\begin{document}

\maketitle

\section*{Note on this draft}
\label{sec:titles}
This is a structured working draft prepared for the \textit{Journal of Regional Analysis \& Policy}. It was
written to the requested outline in full, with the argument, framing, and language choices worked out, but
with all data-dependent content -- data sources, sample coverage, coefficients, standard errors, sample
sizes, and figures -- initially left as clearly marked placeholders (in brackets, e.g., [INSERT ESTIMATE])
because those had not yet been supplied. \textbf{Section~\ref{sec:eventstudy_diD} (event-study extension), Section~\ref{sec:results}'s Section~6.1
(border proximity and revenue per dispensary, including the market-age dynamics), Section~\ref{sec:entry_hazard}
(entry-timing hazard models, including the by-rank, mechanism/placebo, and entry-threshold extensions),
Section~\ref{sec:model_timing} (the entry-timing model), Section~\ref{sec:institutions} (institutional
background for CO/WA/IL/MA/OH), Section~\ref{sec:lit_sequential} and Appendix~\ref{app:literature} (the
sequential-entry literature review), and Section~\ref{sec:discussion_timing} (discussion of the entry-timing
result's limitations and contribution level) now report real, verified content} from the project's
Colorado/Washington county-revenue panel and nine-state entry-density universe, described in full in those
sections' own notes. Every other bracketed item elsewhere in Results and Robustness -- the population-exposure
dose-response, placebo-sector tests, and most of Section~7 -- has not yet been supplied and must still be
filled in from the actual analysis before submission; nothing there should be read as a real reported number. Citations are limited to studies
verified during drafting; any citation the author still needs to confirm is marked ``[citation to
verify].''

\paragraph{Possible titles.}
\begin{enumerate}[label=(\arabic*)]
  \item Prohibition Next Door: Distance, Population Exposure, and Dispensary Revenue at U.S. Cannabis-Market Borders
  \item Captive Demand Across State Lines: Dispensary Revenue and the Border Effects of Partial Marijuana Prohibition
  \item How Far Does Prohibition Travel? Population Exposure and the Cross-Border Cannabis Revenue Premium
  \item Cross-Border Cannabis Demand and Dispensary Revenue: National Evidence from U.S. State Borders
  \item The Revenue Premium at the Border: State Cannabis Policy Gaps and Dispensary Sales in the United States
\end{enumerate}

\paragraph{AI-assistance disclosure.} In accordance with the journal's policy on generative artificial
intelligence, the authors disclose that an AI writing tool (Claude, Anthropic) was used to help structure
and draft portions of this manuscript's prose, including the literature synthesis and exposition of the
empirical strategy, under the authors' direction and subject to the authors' review. All data construction,
empirical analysis, and interpretive conclusions are the responsibility of the authors, and all cited
findings were verified against their original sources before inclusion. [ADAPT OR REMOVE TO MATCH FINAL
JOURNAL DISCLOSURE REQUIREMENTS AT SUBMISSION.]

\newpage

\begin{abstract}
\noindent Recreational-marijuana legalization has proceeded unevenly across U.S. states, leaving many legal
markets directly adjacent to states where recreational sales remain prohibited. This paper asks whether
that policy gap is reflected in dispensary revenue: do counties in legal-market states that are closer to a
nonlegal-state border generate more revenue per active dispensary than counties farther from such a border?
Using a national sample of legal-state counties, we measure each county's distance to the nearest nonlegal
state and the population living on the nonlegal side within reach of nearby dispensaries. We find that
revenue per active dispensary is higher in counties closer to a nonlegal border, and that this relationship
strengthens as the nearby nonlegal-side population grows, consistent with a larger potential cross-border
consumer market. Placebo tests using non-cannabis retail sectors do not show a comparable border pattern,
which is consistent with the result reflecting cannabis-specific demand rather than general border
commerce. Building on evidence that consumers cross state lines to purchase cannabis, we extend this
literature by showing that the resulting dispensary revenue premium varies systematically and continuously
with the size of the population exposed to prohibition nearby, across multiple state borders rather than a
single policy transition. [FINALIZE ABSTRACT LENGTH TO 150--200 WORDS AND INSERT KEY MAGNITUDES ONCE
RESULTS ARE AVAILABLE.]
\end{abstract}

\newpage

\section{Introduction}

Recreational marijuana is legal in a growing number of U.S. states but remains prohibited in most others,
producing sharp policy discontinuities at many state borders. A resident of a nonlegal state can often
reach a licensed dispensary in a neighboring legal state by driving a short distance, while a resident who
stays home cannot purchase recreational cannabis legally at all. This kind of geographic policy
discontinuity is a familiar setting in public economics: when a good is taxed, regulated, or prohibited
differently on the two sides of a border, consumers close to that border have an incentive to cross it, and
that incentive should be strongest where travel costs are low and the population able to cross is large.
[INSERT ONE OR TWO SENTENCES ON THE CURRENT NUMBER OF STATES WITH OPERATING RECREATIONAL MARKETS AND THE
COUNT/SHARE STILL PROHIBITING RECREATIONAL SALES AS OF THE SAMPLE PERIOD.]

Understanding whether this incentive shows up in dispensary revenue matters for several audiences.
State governments that have legalized cannabis face policy questions about licensing, local zoning, and tax
design that depend on how much of dispensary revenue reflects local demand versus demand drawn in from
neighboring, still-prohibited states. State governments that have not legalized cannabis face a fiscal
externality: potential tax revenue and consumption activity that would otherwise occur, or be taxed, within
their own borders instead accrues to a neighboring state. Local governments in border counties face
practical questions about licensing capacity, traffic, and enforcement that depend on whether their
dispensaries serve a substantially larger effective market than their own resident population would imply.
More broadly, the setting offers a way to study how far the economic effects of a state-level policy
discontinuity extend, and how those effects scale with the size of the population on the restricted side of
the border.

This paper asks a direct empirical question: are counties in legal-market states that are closer to the
border of a nonlegal state associated with higher cannabis revenue per active dispensary than counties
farther from such a border? We construct a national sample of counties in states with an operating
recreational-marijuana retail market, measure each county's distance to the nearest state without such a
market, and measure the population living on the nonlegal side of that border within reach of nearby
dispensaries. We find that revenue per active dispensary is higher, and that this relationship is stronger
where the nearby nonlegal-side population is larger, producing a dose-response pattern in which the revenue
premium scales with the size of the potential cross-border consumer market. We report revenue per active
dispensary rather than profit throughout, since the underlying data do not include dispensary operating
costs. To assess whether this pattern is specific to cannabis rather than a general feature of economic
activity near state borders, we compare it against the same specification applied to non-cannabis retail
sectors and do not find a comparable border-proximity pattern in those placebo sectors.

The closest existing work is \citet{HansenMillerWeber2020}, who use Washington's complete cannabis
transaction data and Oregon's 2015 opening of its own recreational retail market to show that Washington
dispensaries near the Oregon border experienced a substantial, immediate decline in sales once Oregon
consumers gained a legal in-state alternative -- direct evidence that consumers had been crossing the
border to buy cannabis. \citet{HaoCowan2020} compare counties directly bordering legal states with counties
farther inside neighboring nonlegal states, but their outcome is marijuana-related arrests rather than
dispensary revenue. Both studies establish, in different ways and in different settings, that state
cannabis policy gaps generate cross-border behavior. This paper does not claim to be the first to document
that consumers cross state borders to buy cannabis; that finding is already well established. Its
contribution is to examine whether the resulting demand is systematically reflected in dispensary revenue
across many legal-nonlegal borders at once, rather than around a single border or a single policy
transition, and to show that the magnitude of the border revenue premium varies continuously with the size
of the population exposed to prohibition nearby, not only with distance. Existing studies establish that
consumers cross state borders to purchase cannabis; we extend this evidence by showing that the resulting
dispensary revenue premium varies systematically with the size of the nearby population living in nonlegal
states, and that this pattern does not appear in placebo retail sectors unrelated to cannabis policy.

The remainder of the paper proceeds as follows. Section~2 reviews the relevant literature on cross-border
shopping, cannabis legalization spillovers, cannabis retail market structure, and substitute goods.
Section~3 develops a simple conceptual framework linking dispensary revenue to distance, nonlegal-side
population, travel cost, legal access, and local supply restrictions. Section~4 describes the data.
Section~5 presents the empirical strategy. Section~6 reports the main results. Section~7 reports robustness
checks. Section~8 discusses policy implications and limitations, and Section~9 concludes.

\section{Literature Review}

\subsection{Cross-Border Shopping and Tax Avoidance}

A long tradition in public economics studies how consumers respond to geographic differences in taxation or
regulation by shopping across a nearby border. \citet{Lovenheim2008} shows that a substantial share of
cigarette consumers in high-tax states purchase cigarettes in nearby lower-tax jurisdictions, and that
demand elasticities with respect to a consumer's home-state price vary systematically with distance to the
nearest lower-price border -- between roughly 13 and 25 percent of consumers in his sample purchase
cigarettes outside their home state. This literature establishes two ideas this paper builds on directly:
first, that the relevant \emph{price} or \emph{legal} differential a consumer faces is defined relative to
the nearest more permissive jurisdiction, not simply their state of residence; and second, that the
intensity of cross-border activity should fall off with distance from the border. [INSERT ADDITIONAL
CROSS-BORDER TAX AVOIDANCE CITATIONS AS APPROPRIATE, E.G. ON ALCOHOL OR GASOLINE TAX BORDER EFFECTS --
CITATION TO VERIFY.]

\subsection{Recreational-Marijuana Legalization and Spillovers}

\citet{HansenMillerWeber2020} provide the most direct existing evidence that consumers cross state borders
specifically to purchase cannabis. Using the near-universe of Washington State's legal cannabis
transactions and the October 2015 opening of Oregon's own recreational retail market as a natural
experiment, they show that Washington retailers near the Oregon border lost a substantial share of their
sales once Oregon consumers no longer needed to cross the border, and that this cross-border diversion cost
Washington millions of dollars in foregone tax revenue. Their design isolates a single border and a single
policy transition -- the opening of Oregon's market -- and shows convincingly that cross-border cannabis
commerce is real and economically meaningful.

\citet{HaoCowan2020} take a complementary approach, comparing marijuana-related arrests in counties that
directly border a legal state against counties located farther inside neighboring nonlegal states. Their
outcome variable, arrests, captures a different margin of the same underlying phenomenon: rather than
measuring how much legal revenue border proximity generates on the legal side, they measure how legal
availability nearby affects enforcement outcomes on the nonlegal side. Together, \citet{HansenMillerWeber2020}
and \citet{HaoCowan2020} establish that legal-nonlegal cannabis policy borders generate real behavioral
responses on both sides. [INSERT ADDITIONAL SPILLOVER CITATIONS ON CRIME, TRAFFIC SAFETY, OR YOUTH USE AS
APPROPRIATE -- CITATION TO VERIFY.]

This paper's contribution relative to this pair of studies is not to show that cross-border cannabis
demand exists -- that is already established -- but to examine whether it is systematically reflected in
dispensary revenue across many legal-nonlegal borders simultaneously, and whether the size of that revenue
response scales with a continuous measure of the nonlegal-side population exposed to it, rather than being
documented at a single border around a single legalization event.

\subsection{Cannabis Retail Markets, Licensing, Entry, and Market Power}

A separate strand of work studies the structure of legal cannabis retail markets themselves.
\citet{HollenbeckUetake2021} show that Washington State's restrictive licensing regime, which caps the
number of retail licenses available, produced substantial market power among incumbent retailers, with
consequences for how tax incidence and pricing respond to policy changes. This work is directly relevant to
interpreting a revenue-per-dispensary outcome: if the number of licensed dispensaries in a market is capped
by regulation rather than determined by free entry, then a shock to local demand -- such as demand drawn in
from a nearby nonlegal state -- need not translate into more dispensaries opening. Instead, it can appear
entirely as higher revenue per existing dispensary, which is the outcome this paper studies. The classical
economic logic connecting market size, demand, and the number of firms a market can support in equilibrium
is developed in \citet{BresnahanReiss1991}, whose framework for analyzing entry in concentrated markets
underlies the distinction this paper draws between the number of active dispensaries and revenue per
dispensary as separate margins of market response. Washington's dispensary licenses were themselves
allocated through a lottery in oversubscribed jurisdictions
\citep{WSLCB2014LotteryPR,DongTyndall2024,ThomasTian2021}, and Illinois, Massachusetts, and Ohio each impose
their own distinct entry frictions -- merit-based scoring with a lottery for tied applicants
\citep{IDFPRAdultUseCannabis}, a locally negotiated host community agreement \citep{MACCCLicensing}, and a
staged conversion of existing medical dispensaries followed by a separate capped round for new entrants
\citep{OhioDCCFAQ}, respectively -- discussed further in Section~\ref{sec:data}.

\subsection{Sequential Entry and the Economics of the Second Entrant}
\label{sec:lit_sequential}

A separate empirical result in this paper concerns how quickly a second dispensary follows the first in a
given county, rather than the level of revenue an existing dispensary earns. This connects to a distinct,
older strand of the entry literature, which we summarize briefly here and treat at greater length in
Appendix~\ref{app:literature}. \citet{BresnahanReiss1990} and \citet{BresnahanReiss1991} establish, using
market-size thresholds across several concentrated retail and professional industries, that competitive
conduct changes most sharply with the arrival of a market's second or third firm, and that additional entry
beyond that point has little further effect -- a market-size, cross-sectional version of the same basic
idea our timing result speaks to. \citet{Berry1992}, \citet{Mazzeo2002}, and \citet{Seim2006} extend this
tradition to settings with firm heterogeneity and endogenous product-type choice. \citet{ToivanenWaterson2005}
is the closest dynamic analogue: using a fast-food duopoly panel, they show a rival's existing presence is a
strong predictor of a chain's own subsequent entry decision. \citet{IgamiYang2016} model sequential,
preemption-motivated multi-store entry directly. None of this literature has, to our knowledge, been applied
to cannabis dispensary markets, and none of it isolates a distinctively \emph{second}-entrant effect using a
timing or hazard-based design of the kind we use in Section~\ref{sec:entry_hazard}; our contribution here is
best described as applying a known regularity from the general entry literature, using a different empirical
method, to a new and highly regulated retail setting, rather than as identifying a new mechanism.

\subsection{Alcohol and Other Possible Substitute Goods}

Because cannabis legalization can affect the demand for related goods, cannabis-specific claims require
some caution when a comparison sector is itself plausibly affected by cannabis policy. \citet{BaggioChongKwon2020}
use retail scanner data on alcohol purchases across U.S. counties and a border-pair design to show that
counties in states with medical marijuana laws reduced monthly alcohol sales relative to nearby counties in
states without such laws, consistent with cannabis and alcohol acting as substitutes in at least some
contexts. This finding has a direct methodological implication for this paper: alcohol retail is not a
clean placebo sector for a cannabis border-revenue test, because marijuana-policy differences across a
border could independently raise or lower alcohol revenue through substitution or complementarity, rather
than leaving it unaffected. We include alcohol among the comparison sectors reported in the results because
it is informative and because it appears in related work, but we treat any alcohol-sector placebo result
with more caution than results from sectors -- such as general retail trade unrelated to cannabis or
alcohol -- that are not plausibly affected by marijuana policy through a demand-substitution channel. [INSERT
ADDITIONAL CITATIONS ON CANNABIS-ALCOHOL OR CANNABIS-TOBACCO SUBSTITUTION AS APPROPRIATE -- CITATION TO
VERIFY.]

\section{Conceptual Framework}

We motivate the empirical analysis with a simple, informal framework linking a dispensary's revenue to the
demand it can draw from both sides of a nearby state border and to the supply-side constraints it operates
under. Consider a county $i$ in a legal-market state, located a distance $d_i$ from the nearest county in a
neighboring state where recreational cannabis remains prohibited. Residents of that nonlegal-side county
face a binary choice: forgo legal recreational cannabis, obtain it illegally, or travel to a nearby legal
dispensary. The number of residents for whom the last option is attractive depends on the size of the
nonlegal-side population within a feasible travel radius, denoted $P_i$, and on the travel cost of reaching
county $i$'s dispensaries, which rises with distance $d_i$ and with any adjustment for road access or
travel time rather than straight-line distance alone. It also depends on legal access: only adults who meet
the minimum purchase age and can present valid identification may buy legally, so $P_i$ should in principle
be measured as the eligible adult population rather than total population, to the extent that data allow
this distinction.

The demand a county $i$'s dispensaries can draw from the nonlegal side, $D_i$, is therefore increasing in
$P_i$ and decreasing in $d_i$ and in the travel cost associated with $d_i$:
\begin{equation}
D_i = D(P_i, d_i), \qquad \frac{\partial D_i}{\partial P_i} > 0, \qquad \frac{\partial D_i}{\partial d_i} < 0.
\end{equation}
This additional demand is added to whatever local, in-state demand county $i$'s dispensaries already serve.
Total realized demand is met by $N_i$ active dispensaries in county $i$. If $N_i$ adjusts freely to meet
demand -- as in a standard model of entry where the number of firms a market supports rises with market
size, following the entry-and-market-size logic in \citet{BresnahanReiss1991} -- then a larger $D_i$ should
appear partly or wholly as more dispensaries rather than as higher revenue per dispensary. But where entry
is constrained by state or local licensing rules, occupancy or zoning restrictions, or outright local bans,
as documented for at least one major legal-cannabis state by \citet{HollenbeckUetake2021}, $N_i$ cannot
adjust freely. In that case, additional cross-border demand should appear as higher revenue per existing
dispensary, $R_i / N_i$, rather than as a larger $N_i$:
\begin{equation}
\frac{R_i}{N_i} = f\big(D_i(P_i, d_i),\ \text{local demand}_i,\ N_i\big), \qquad
\frac{\partial (R_i/N_i)}{\partial P_i} > 0 \ \text{when } N_i \text{ is supply-constrained}.
\end{equation}
This framework yields three testable implications that map directly onto the empirical strategy below.
First, revenue per dispensary should be higher in counties closer to a nonlegal border, holding local
demand fixed, because $d_i$ is smaller. Second, this border-proximity effect should be stronger where the
nonlegal-side population $P_i$ is larger, producing an interaction, or dose-response, pattern between
distance and nonlegal-side population. Third, because the mechanism is specific to cannabis policy rather
than a general feature of economic activity near state lines, an analogous specification applied to
non-cannabis retail sectors -- for which no comparable legal discontinuity exists at the same border --
should not show the same pattern, providing a placebo test of the interpretation. The framework also makes
clear why the number of active dispensaries, entry and exit, and total county-level cannabis revenue are
each informative complementary outcomes: they help distinguish whether an observed border effect operates
through supply-constrained pricing, through the extensive margin of dispensary entry, or through some
combination of the two, which is why Section~6 and Section~7 report results for each of these outcomes
alongside the primary revenue-per-dispensary measure.

\subsection{Entry Timing and the Second Dispensary}
\label{sec:model_timing}

The framework above is static: it describes the demand a county can support once its dispensary sector has
settled at $N_i$ stores, but it is silent on \emph{when} a county reaches a given $N_i$. Section~\ref{sec:entry_hazard}
reports a result the static framework alone cannot speak to: distance to a nonlegal border predicts how
quickly a county's \emph{second} dispensary follows its first, but not how quickly the first dispensary
follows a state's own retail launch, and the effect fades by the third and fourth entrant. We extend the
framework with a minimal timing dimension to make sense of this pattern, following the spirit of sequential
entry-under-uncertainty models such as \citet{PrescottVisscher1977}.

A regulatory or capital-raising delay $F_i$ -- licensing review, local approval, buildout -- must be cleared
before any dispensary can open, and we treat this delay as largely orthogonal to a county's realized demand:
a prospective entrant must clear the same administrative queue whether or not the county turns out to have
much captive cross-border demand. This is consistent with Table~\ref{tab:entry_hazard_by_rank}'s finding
that a county's own retail-launch-to-first-dispensary time shows no relationship with $d_i$. What a first
mover's entry \emph{does} do, once open, is publicly reveal information about realized demand -- sales
volume, wait times, local press coverage -- that was not observable beforehand. Local population $P_i$ is
known in advance from public data and does not need to be revealed; whether nonlegal-side residents actually
cross the border in practice is the part of demand that is genuinely uncertain ex ante. A potential second
entrant, deciding whether to enter after observing the first store's success, therefore places more weight
on the border-proximity channel specifically than the first entrant's ex ante decision could, because the
first entrant had to price that uncertainty at its average value across otherwise similar counties, while
the second entrant can condition on a realized signal. This generates the central asymmetry in
Table~\ref{tab:entry_hazard_by_rank}: distance should matter more for the second entrant's timing than the
first's. As a county's dispensary count approaches whatever level its now-revealed demand $D_i$ supports,
further entry is governed by ordinary population growth rather than the one-time resolution of uncertainty
about cross-border demand, which is consistent with the effect fading by the third and fourth entrant.

This account also offers one candidate explanation, though not the only plausible one, for why
Table~\ref{tab:entry_hazard_mechanism} finds that the \emph{size} of the nonlegal-side population does not
predict second-entrant timing once a county is within reach of some nonlegal border, while proximity itself
does: if crossing the border to shop is a discrete decision governed by a travel-cost threshold rather than
a continuously demand-elastic one, a second entrant may only need to learn \emph{whether} meaningful
cross-border traffic exists at all, not precisely how large the population behind it is, for a given county
to look profitable. An equally plausible alternative is that entrants use border proximity as a coarse
heuristic rather than a calibrated demand estimate. Our data cannot distinguish between these two
microfoundations, and we do not take a stand on which is correct; both predict the same qualitative pattern
we observe; and we note explicitly that Table~\ref{tab:entry_threshold} does not corroborate a further
implication one might expect from the demand-supplement story -- that border counties should also reach the
second-store population threshold sooner in levels, not just faster in time -- which we discuss as a real
limitation in Section~\ref{sec:results} rather than set aside.

Formally, let $\hat{S}_i$ denote the market size a second entrant infers after observing the first
dispensary's revealed performance, updated from a common prior $S_i^0$ using a signal that is more
informative about the border-demand component the closer $i$ is to a nonlegal county:
\begin{equation}
\hat{S}_i = S_i^0 + \lambda(d_i) \cdot \varepsilon_i, \qquad \lambda'(d_i) < 0,
\end{equation}
where $\varepsilon_i$ is the revealed surprise in realized demand and $\lambda(d_i)$ is the weight placed on
that surprise, decreasing in distance because a nearby nonlegal border is the main source of demand
uncertainty the first entrant resolves. The second entrant opens once $\pi_2(\hat{S}_i) \geq F$. This
delivers five comparative statics, the first four of which match Tables~\ref{tab:entry_hazard_by_rank}
through \ref{tab:entry_threshold} and the fifth of which they do not:
\begin{enumerate}
  \item First-entry timing does not depend systematically on $d_i$, because $F_i$ dominates and $\hat{S}_i$
        has not yet been revealed to anyone.
  \item Second-entry timing is decreasing in $d_i$, because $\lambda(d_i)$ is decreasing in $d_i$.
  \item If $\lambda(d_i)$ falls off sharply beyond some reach $\bar{d}$ (a travel-cost threshold) rather than
        smoothly with the exact size of the population behind it, second-entry timing should respond to
        proximity more than to population size conditional on being within reach.
  \item The effect should fade by the third and fourth entrant as $N_i$ approaches its now-revealed,
        fundamentals-implied level.
  \item \emph{(Not supported by our data.)} If cross-border demand genuinely raises a county's effective
        market size, border counties should also reach the population level needed to support a second
        store sooner than interior counties, in the same way they already do for the first store. We do not
        find this pattern (Table~\ref{tab:entry_threshold}) and report it as an open anomaly rather than
        omit it.
\end{enumerate}

\section{Data}
\label{sec:data}

[THIS SECTION IS A TEMPLATE. EVERY BRACKETED ITEM MUST BE REPLACED WITH THE ACTUAL DATA DESCRIPTION BEFORE
SUBMISSION.]

\subsection{Data Sources}
The primary outcome data come from [INSERT DATA SOURCE(S), E.G. STATE DEPARTMENT OF REVENUE OR CANNABIS
CONTROL AGENCY ADMINISTRATIVE RECORDS, BY STATE]. Dispensary license and location data come from [INSERT
DATA SOURCE(S), E.G. STATE LICENSING REGISTRIES, GEOCODED TO COUNTY OR CENSUS TRACT]. Population data come
from [INSERT SOURCE, E.G. U.S. CENSUS BUREAU POPULATION ESTIMATES]. Comparison-sector (placebo) revenue
data come from [INSERT SOURCE, E.G. STATE DEPARTMENT OF REVENUE RETAIL SALES-TAX FILINGS BY INDUSTRY AND
COUNTY].

\subsection{Sample Period and Coverage}
The sample covers [INSERT YEARS COVERED] and includes counties in [INSERT LIST OF STATES INCLUDED, WITH
RETAIL-LEGALIZATION DATES]. States are included if [INSERT INCLUSION CRITERION, E.G. COUNTY-LEVEL REVENUE
DATA WERE OBTAINABLE FROM THE STATE REGULATOR]. [INSERT ANY STATES CONSIDERED BUT EXCLUDED FOR DATA
AVAILABILITY REASONS, AND WHY.]

\subsection{Institutional Background: Licensing Regimes in the Entry-Timing Sample}
\label{sec:institutions}

The entry-timing analysis in Section~\ref{sec:entry_hazard} uses a broader, nine-state sample (California,
Colorado, Illinois, Massachusetts, Maine, New York, Ohio, Oregon, and Washington) than the two-state
revenue-per-store analysis, and the states doing most of the statistical work for that result -- Illinois,
Massachusetts, and Ohio -- have licensing regimes distinct from Colorado and Washington's. We summarize the
features relevant to entry timing and sequencing here, deferring general institutional history to
Appendix~\ref{app:literature}.

\textbf{Washington} capped the statewide number of retail licenses at 334 from the outset and allocated them
by lottery in any jurisdiction where applications exceeded the local quota \citep{WSLCB2014LotteryPR}, a
design also used as a source of exogenous variation in \citet{DongTyndall2024} and \citet{ThomasTian2021}.
Because many of a jurisdiction's licenses could be drawn in the same lottery round, a county's first and
second dispensaries need not have opened sequentially in response to observed demand; they may instead have
opened close together because both winners cleared the same administrative step at the same time. This is a
plausible reason Washington shows no second-entrant timing effect in Table~\ref{tab:entry_hazard_mechanism}
even though its licensing system otherwise resembles the capped regimes that do show the effect.

\textbf{Colorado} imposes no statewide numerical cap; Amendment 64 instead gives local governments full
authority to prohibit retail marijuana outright or to regulate its number, place, and manner, and a majority
of Colorado municipalities have exercised this option to ban it entirely. Where permitted, entry is a more
continuous, decentralized process without a single statewide allocation event.

\textbf{Illinois} initially awarded 75 adult-use dispensary licenses using a merit-based points system (up to
250 points, including social-equity bonus points), with ties among the highest scorers broken by lottery; a
further 110 licenses were awarded through two additional lotteries restricted to social-equity applicants
\citep{IDFPRAdultUseCannabis}.

\textbf{Massachusetts} requires every applicant to negotiate and sign a Host Community Agreement with its
municipality, in addition to state licensing and local zoning approval \citep{MACCCLicensing} -- a
bilateral, locally negotiated step with no fixed timeline, which could plausibly resolve faster for a second
applicant once a municipality has already negotiated one such agreement.

\textbf{Ohio} legalized adult-use retail in November 2023 (Issue 2) and initially converted each existing
medical dispensary directly into an adult-use licensee, with dispensaries not commonly owned by a cultivator
or processor receiving an additional license; a separate, capped round of up to 50 new adult-use dispensary
licenses for applicants without an existing medical license followed later \citep{OhioDCCFAQ}.

These institutional differences mean the pooled entry-timing sample combines counties whose first and
second entrants arose from at least four qualitatively different administrative processes -- a statewide
lottery, a fully decentralized local-option system, a merit-score-plus-lottery system, and a medical-to-
adult-use conversion followed by a separate capped round. We are not able to fully separate which of these
regimes' features drive the pooled result beyond the state-by-state breakdown already reported in
Table~\ref{tab:entry_hazard_mechanism}, and we flag this as a genuine limitation on causal interpretation in
Section~\ref{sec:results}.

\subsection{Geographic Unit}
The unit of observation is [INSERT GEOGRAPHIC UNIT, E.G. COUNTY-YEAR OR COUNTY-QUARTER]. [INSERT NUMBER OF
COUNTIES AND TOTAL OBSERVATIONS IN THE FINAL SAMPLE.]

\subsection{Revenue Definition}
The primary outcome is log revenue per active dispensary, defined as [INSERT PRECISE DEFINITION: TOTAL
REPORTED TAXABLE OR GROSS CANNABIS SALES IN A COUNTY-PERIOD, DIVIDED BY THE NUMBER OF ACTIVE DISPENSARIES IN
THAT COUNTY-PERIOD]. We refer to this measure as revenue per dispensary, not profit, throughout the paper,
because the underlying administrative data do not include dispensary operating costs, rents, wages, or
other expenses. [INSERT WHETHER MEDICAL AND RECREATIONAL SALES ARE COMBINED OR REPORTED SEPARATELY, AND
WHICH IS USED AS THE PRIMARY OUTCOME.]

\subsection{Definition of an Active Dispensary}
A dispensary is classified as active in a given period if [INSERT DEFINITION, E.G. ITS LICENSE STATUS
INDICATES OPEN/OPERATING AS OF THAT PERIOD'S REPORTING CUTOFF, OR IT REPORTS NONZERO SALES]. County-periods
with zero active dispensaries are dropped from the revenue-per-dispensary sample, since the outcome is
undefined; [INSERT NUMBER OF OBSERVATIONS DROPPED FOR THIS REASON].

\subsection{Distance Measure}
The primary distance measure is [INSERT DEFINITION, E.G. STRAIGHT-LINE (GREAT-CIRCLE) DISTANCE FROM THE
COUNTY POPULATION CENTROID, OR FROM EACH DISPENSARY'S GEOCODED LOCATION, TO THE NEAREST COUNTY IN A STATE
WITHOUT AN OPERATING RECREATIONAL RETAIL MARKET]. [STATE WHETHER A DRIVING-TIME OR ROAD-NETWORK DISTANCE
MEASURE IS ALSO CONSTRUCTED, AND IF SO, ITS SOURCE, E.G. A ROUTING API OR NETWORK-DISTANCE SOFTWARE.] The
main specification uses a border threshold of [INSERT THRESHOLD, E.G. A NUMBER OF MILES], with robustness
checks at alternative thresholds reported in Section~7.

\subsection{Population-Exposure Measure}
The nearby nonlegal-side population measure is constructed as [INSERT DEFINITION, E.G. THE SUM OF POPULATION
IN ALL NONLEGAL-STATE COUNTIES WITHIN A GIVEN RADIUS OF EACH LEGAL-STATE COUNTY OR DISPENSARY, USING A
RADIUS QUERY AGAINST COUNTY POPULATION CENTROIDS]. [INSERT WHETHER THIS IS RESTRICTED TO THE ADULT
POPULATION OF LEGAL PURCHASE AGE, AND THE RADIUS OR RADII USED.]

\subsection{Comparison-Sector (Placebo) Data}
Comparison-sector revenue data cover [INSERT SECTORS INCLUDED, E.G. GENERAL RETAIL TRADE, GROCERY,
RESTAURANTS, AND ALCOHOL RETAIL], drawn from [INSERT SOURCE] at the same geographic and time granularity as
the cannabis revenue data, where available. [INSERT ANY COVERAGE LIMITATIONS, E.G. DISCLOSURE SUPPRESSION
FOR SMALL COUNTIES OR INDUSTRIES.] As discussed in Section~2.4, alcohol retail is included but interpreted
cautiously given evidence that cannabis and alcohol demand may not be independent
\citep{BaggioChongKwon2020}; sectors with no plausible direct link to cannabis policy provide the primary
placebo evidence.

\subsection{Summary Statistics}
[INSERT SUMMARY STATISTICS TABLE: NUMBER OF COUNTIES, NUMBER OF STATES, NUMBER OF BORDER VS. INTERIOR
COUNTIES AT THE MAIN THRESHOLD, MEAN AND STANDARD DEVIATION OF REVENUE PER DISPENSARY, NUMBER OF ACTIVE
DISPENSARIES, DISTANCE TO NEAREST NONLEGAL BORDER, AND NONLEGAL-SIDE POPULATION EXPOSURE.]

\section{Empirical Strategy}
\label{sec:empirical}

\subsection{Main Border-Proximity Specification}
The main specification relates log revenue per active dispensary in county $i$, state $s$, period $t$ to
that county's proximity to the nearest nonlegal-state border:
\begin{equation}
\ln\left(\frac{\text{Revenue}_{ist}}{\text{Dispensaries}_{ist}}\right) =
\alpha + \beta_1\, \text{Border}_i + \beta_2\, \ln(\text{Population}_{ist}) + X_{ist}'\gamma +
\theta_t + \delta_s + \varepsilon_{ist},
\end{equation}
where $\text{Border}_i$ is either a binary indicator equal to one if county $i$ is within [INSERT THRESHOLD]
of a nonlegal-state border, or the continuous log distance to that border (with sign reversed for
interpretation, or entered so that a negative coefficient indicates revenue rising as distance falls,
consistent with the framework in Section~3); $X_{ist}$ is a vector of time-varying county controls [INSERT
CONTROLS, E.G. LOG LAND AREA, URBANICITY, MEDIAN INCOME, LOCAL TAX RATE, YEARS SINCE THE COUNTY'S OWN STATE
LEGALIZED RETAIL SALES]; $\theta_t$ are period fixed effects; and $\delta_s$ are state fixed effects, which
absorb any systematic differences in cannabis policy, taxation, or market conditions across states that are
unrelated to border proximity. Because border status is a fixed geographic attribute of a county, $\beta_1$
is identified from the cross-sectional and within-state contrast between border and interior counties,
conditional on the controls, not from counties switching status over time.

\subsection{Dose-Response Specification}
To test whether the border-proximity relationship strengthens with the size of the nearby nonlegal-side
population, we augment the main specification with the population-exposure measure and its interaction with
border proximity:
\begin{equation}
\ln\left(\frac{\text{Revenue}_{ist}}{\text{Dispensaries}_{ist}}\right) =
\alpha + \beta_1\, \text{Border}_i + \beta_2\, \ln(\text{NonlegalPop}_i) +
\beta_3\, \big(\text{Border}_i \times \ln(\text{NonlegalPop}_i)\big) +
X_{ist}'\gamma + \theta_t + \delta_s + \varepsilon_{ist}.
\end{equation}
A positive $\beta_3$ (or, in the fully continuous version described next, a positive coefficient on the
interaction between log distance and log nonlegal-side population with the appropriate sign convention)
would indicate that the revenue premium associated with border proximity is larger where more people live
on the nonlegal side within reach, consistent with a larger potential cross-border consumer market. We also
estimate a fully continuous version, replacing the binary border indicator with log distance to the nearest
nonlegal border interacted with log nonlegal-side population, to characterize the dose-response relationship
without relying on a single distance threshold:
\begin{equation}
\ln\left(\frac{\text{Revenue}_{ist}}{\text{Dispensaries}_{ist}}\right) =
\alpha + \beta_1\, \ln(d_i) + \beta_2\, \ln(\text{NonlegalPop}_i) +
\beta_3\, \big(\ln(d_i) \times \ln(\text{NonlegalPop}_i)\big) + X_{ist}'\gamma + \theta_t + \delta_s +
\varepsilon_{ist}.
\end{equation}
This interpretation is offered as mechanism evidence consistent with a cross-border demand channel, not as
an independent causal test of that channel on its own: a positive interaction is consistent with the
demand-exposure story in Section~3, but confirming it directly would require information on where
dispensary customers actually reside, which the revenue data used here do not provide.

\subsection{Placebo-Sector Specifications}
To assess whether an estimated border-proximity effect is specific to cannabis, we re-estimate the main and
dose-response specifications with the cannabis revenue-per-dispensary outcome replaced by the analogous
outcome for each comparison sector $k$:
\begin{equation}
\ln\left(\frac{\text{Revenue}^k_{ist}}{\text{Establishments}^k_{ist}}\right) =
\alpha + \beta_1\, \text{Border}_i + \beta_2\, \ln(\text{Population}_{ist}) + X_{ist}'\gamma +
\theta_t + \delta_s + \varepsilon_{ist},
\end{equation}
for sectors $k \in \{$general retail, [INSERT ADDITIONAL SECTORS], alcohol$\}$. Because marijuana
legalization may itself affect alcohol demand through substitution or complementarity
\citep{BaggioChongKwon2020}, we treat any alcohol-sector result as a secondary check and give primary
placebo weight to sectors with no plausible direct connection to cannabis policy. A precise, insignificant
$\beta_1$ in those sectors, alongside a positive and significant $\beta_1$ for cannabis, is the pattern
consistent with a cannabis-specific rather than a general border-commerce interpretation.

\subsection{Event-Study and Difference-in-Differences Extension}
\label{sec:eventstudy_diD}
Two such policy-timing events exist in the underlying county-revenue panel and are usable here. First,
\citet{HansenMillerWeber2020}'s own setting: Oregon opened its retail market on October 1, 2015, more than
a year after Washington's July 2014 launch, so Washington counties bordering Oregon held a temporary
captive-demand advantage that should have switched off at that date. Second, a policy-timing event
\citeauthor{HansenMillerWeber2020} did not study: New Mexico opened its own retail market on April 1, 2022,
which should erode the advantage of Colorado counties bordering New Mexico relative to Colorado's interior
and relative to Colorado counties that still border a permanently nonlegal state (Utah, Wyoming, Oklahoma,
Nebraska, Kansas). We estimate
\begin{equation}
\ln\left(\frac{\text{Revenue}_{ist}}{\text{Dispensaries}_{ist}}\right) =
\alpha + \beta\, \big(\text{Border}_i \times \text{Post}_t\big) +
X_{ist}'\gamma + \theta_t + \varepsilon_{ist},
\end{equation}
where $\text{Border}_i$ is the relevant neighbor-specific border indicator (Oregon or New Mexico, at the
paper's 70-mile threshold) and $\text{Post}_t$ indicates the period after the neighboring state's own
retail market opened, and separately estimate the full event-time coefficient path
$\sum_{k\neq -1}\beta_k(\text{Border}_i\times\mathbb{1}[\text{event time}=k])$ omitting the year immediately
before the event as the reference category.

The two settings diverge. The Colorado-New Mexico extension works cleanly: with the transition year (2022)
dropped, revenue per store in New Mexico-border counties falls 38.6 percent relative to Colorado's interior
after New Mexico opens ($p=0.0001$, county-clustered), and the event-time path is positive and often
significant throughout the eight available pre-years, then declines monotonically from the event year
onward to a precise near-zero by three years after. This magnitude is close to
\citeauthor{HansenMillerWeber2020}'s own 36 percent estimate for an entirely different border and
neighboring-state legalization. The Washington-Oregon replication does not confirm this pattern in our own
county-revenue construction: Washington's revenue data are annual (fiscal year, July-June) rather than
\citeauthor{HansenMillerWeber2020}'s monthly administrative records, so fiscal year 2016 straddles the
October 2015 opening and is a genuine transition period rather than a clean post period; even so, the
Oregon-border coefficient is wrong-signed (positive rather than negative) and statistically indistinguishable
from zero in both the including- and excluding-transition-year specifications ($p\approx0.52$-$0.53$,
county-clustered), climbing steadily across the following decade rather than dropping at the 2015 switch.
This does not survive excluding Washington's two largest, most populous Oregon-border counties (Clark and
Thurston), so it is not a single-county artifact; we read it as a limitation of annual, current-dispensary-
count data relative to \citeauthor{HansenMillerWeber2020}'s monthly, store-level records, not as evidence
against their finding, which they establish directly. [INSERT FINAL TABLE NUMBER FOR THE COMBINED
EVENT-STUDY SUMMARY (see \texttt{tables/co\_nm\_eventstudy\_vs\_interior.csv},
\texttt{tables/co\_nm\_eventstudy\_vs\_otherborder.csv}, \texttt{tables/co\_nm\_did\_summary.csv},
\texttt{tables/wa\_or\_did\_summary.csv}, and \texttt{tables/wa\_or\_eventstudy\_path.csv} in this draft's
own \texttt{tables/} folder, and \texttt{figures/fig\_co\_nm\_eventstudy.png} /
\texttt{figures/fig\_wa\_or\_eventstudy.png} for the event-time plots) ONCE THIS SECTION IS INTEGRATED WITH
THE REST OF THE DRAFT'S TABLE NUMBERING.]

\subsection{Inference}
Standard errors are clustered at [INSERT CLUSTERING LEVEL, E.G. THE STATE LEVEL, THE BORDER-PAIR LEVEL, OR
THE COUNTY LEVEL], reflecting the level at which the underlying policy variation -- a state's own
legalization status -- is assigned. [INSERT ROBUSTNESS TO ALTERNATIVE CLUSTERING LEVELS, AND NOTE
EXPLICITLY IF THE NUMBER OF CLUSTERS AT ANY LEVEL IS SMALL ENOUGH TO WARRANT A WILD BOOTSTRAP OR OTHER
SMALL-CLUSTER CORRECTION.]

\subsection{Identifying Assumptions and Threats to Validity}
The main specification is a cross-sectional or panel comparison of border and interior counties within the
same states, not a design with a clean treatment-timing switch, since border status is a fixed geographic
attribute. This has several implications for interpretation. First, border and interior counties may differ
systematically in population density, income, urbanization, tourism activity, road access, retail activity
more broadly, local tax rates, and enforcement intensity, independent of cannabis policy; state and period
fixed effects and the control variables in $X_{ist}$ are intended to absorb the most important of these
differences, but omitted variation at the county level cannot be fully ruled out without a design that
exploits a change over time. Second, dispensaries may locate deliberately near large cities, highways,
tourist destinations, or state borders for reasons unrelated to prohibition-driven demand, which could
generate a spurious correlation between border proximity and revenue if location choice itself responds to
factors correlated with distance to the border; we address this partly through the placebo-sector tests,
which would be expected to show a similar pattern if generic locational advantages, rather than
cannabis-specific demand, were driving the result. Third, revenue per dispensary conflates a genuine demand
response with any effect of licensing restrictions on the number of dispensaries; Section~6 and Section~7
therefore also report results for total county revenue and for the number of active dispensaries, to
clarify which margin is responding. Given these considerations, we describe the main estimates as evidence
that border proximity and nonlegal-side population exposure are associated with higher dispensary revenue,
and as consistent with a cross-border demand interpretation, rather than as strictly causal estimates,
except where the event-study extension in Section~5.4 supports a stronger causal reading around a specific
policy-timing event.

\section{Results}
\label{sec:results}

[THIS SECTION REPORTS ONLY THE QUALITATIVE PATTERN OF THE AUTHORS' FINDINGS AS SUPPLIED. NO COEFFICIENT,
STANDARD ERROR, SAMPLE SIZE, OR $P$-VALUE IN THIS SECTION IS A REAL REPORTED NUMBER; EVERY BRACKETED ITEM
MUST BE FILLED IN FROM THE ACTUAL ESTIMATES BEFORE SUBMISSION.]

\subsection{Border Proximity and Revenue per Dispensary}
Table~\ref{tab:revenue_perstore_border70} reports estimates of the main border-proximity specification, for
the two states with a usable county-year revenue panel and geocoded dispensary counts, Colorado and
Washington (71 counties, $n=706$ county-years with at least one open store). Counties within 70 miles of a
nonlegal-state border have significantly higher log revenue per active dispensary than interior counties of
the same two states: $\hat\beta_1=+0.502$ (HC1 $p=0.031$) with no controls, falling only modestly to
$+0.395$ (county-clustered $p=0.015$) once log(population), year fixed effects, and state fixed effects are
added -- the relationship is robust to that full control sequence, not an artifact of population or common
time trends. In percentage terms this is a roughly 48 percent revenue-per-store premium in border counties.
The result is not driven by a single county: excluding Sedgwick County, Colorado (a small, heavily-touristed
county on the Nebraska border that is a clear geographic outlier in the sample), the coefficient falls only
to $+0.302$ and remains significant ($p=0.040$). It also holds within each state separately: Colorado alone
gives $+0.552$ ($p=0.005$), Washington alone $+0.441$ ($p=0.077$, marginal but directionally consistent).
For comparison, Table~\ref{tab:revenue_border70} reports the identical border/interior specification applied
to revenue \emph{per capita} rather than per dispensary -- the population-adjusted analogue of total county
revenue, since the underlying regulator data report county-aggregate sales rather than a per-transaction
customer count. That specification is a precise null in every column ($\hat\beta_1=+0.046$, county-clustered
$p=0.833$ with full controls; $-0.060$, $p=0.763$ excluding Sedgwick), consistent with the framework in
Section~3: if the number of dispensaries in a county is constrained by state licensing policy rather than
free to expand with demand, a captive-demand effect should appear as higher revenue at each existing store
rather than as higher aggregate per-capita revenue or a larger number of stores -- exactly the divergence
observed between Tables~\ref{tab:revenue_perstore_border70} and \ref{tab:revenue_border70}.

\input{tables/table3_revenue_perstore_border70}

\input{tables/table1_revenue_border70}

\subsubsection{Does the premium hold up, or fade, as a state's legal market matures?}
A cross-sectional border/interior gap could in principle be a transient feature of a market's earliest,
least-competitive years rather than a persistent structural feature. Table~\ref{tab:revenue_perstore_eventstudy}
re-estimates the border-proximity coefficient separately by each state's own market age (years since its
retail launch), controlling for log(population), state fixed effects, and market-age-year fixed effects.
Market-age year 0 is too thin to interpret (only 4 border observations, drawn from both states' earliest
data), but from year 1 onward the coefficient is positive in every single bin -- there is no sign reversal
at any point in either state's market lifecycle -- and grows from marginal or noisy in years 2--4 to
significant and rising from $+0.40$ at year 5 to $+0.59$ at year 9 and beyond.

\input{tables/table4_revenue_perstore_eventstudy}

Figure~\ref{fig:perstore_eventstudy}
plots this coefficient path directly. The premium \emph{strengthens} rather than fades as the market
matures, which we read as more consistent with a binding, roughly time-invariant supply constraint (a state
licensing cap that becomes more binding relative to demand as the overall market grows) than with a
temporary novelty effect that a competitive process would be expected to erode over time.

\begin{figure}[htbp]
\centering
\includegraphics[width=0.9\textwidth]{figures/fig3_revenue_perstore_eventstudy.png}
\caption{Revenue per store: border $\times$ market-age event-study coefficients, with 95\% confidence
intervals (county-clustered standard errors). There is no pre-launch period to plot, since no legal
cannabis revenue exists before a state's own retail market opens, so the path starts at market-age year 0.}
\label{fig:perstore_eventstudy}
\end{figure}

Figure~\ref{fig:perstore_trends} shows the same underlying data a different way, plotting border and
interior counties' revenue-per-store \emph{levels} separately (rather than only their difference) on the
same market-age clock. The two groups start close together in a market's first two years, when the whole
market is still small, and then visibly diverge from year 2 onward, with border counties pulling ahead and
staying ahead through year 9 and beyond. Because neither state has revenue data from before its own retail
launch, this figure cannot establish that the two groups were on parallel trajectories \emph{before}
treatment in the way a true pre-period would; what it shows is that the two groups are not visibly different
in the earliest available data and diverge only afterward, which is the most this data can offer in place of
a formal pre-trend test.

\begin{figure}[htbp]
\centering
\includegraphics[width=0.9\textwidth]{figures/fig4_border_vs_interior_trends_perstore.png}
\caption{Revenue per store: border and interior counties' levels plotted separately (not as a difference),
stacked on a common market-age clock. The two groups start together and diverge from year 2 onward.}
\label{fig:perstore_trends}
\end{figure}

\subsection{Entry Timing: Did Sellers Race to Capture Border Demand?}
\label{sec:entry_hazard}
The revenue results above establish that border demand shows up as higher revenue at existing stores, but
they are silent on a related behavioral question: did sellers \emph{perceive} the border opportunity and
act on it by opening faster, or did entry happen on its own schedule with demand sorting itself out
afterward? This is a different question from the static entry-density/count check discussed in
Section~\ref{sec:data} (which asks how many dispensaries a county ends up with, and finds no border effect),
because it uses each dispensary's own opening date to ask about the \emph{speed} of entry rather than the
eventual count. We estimate Cox proportional-hazards models, stratified by state, over the full nine-state,
479-county entry-density universe described in Section~\ref{sec:data}, for two outcomes: the time from a
county's own state's retail launch to its first dispensary opening, and the time from that first dispensary
to a second one, both right-censored at the data cutoff for counties that have not yet reached the relevant
milestone.

\input{tables/table16_entry_hazard}

The evidence is mixed, and we think the contrast between the two outcomes is itself informative. There is no
evidence that border counties get their \emph{first} dispensary faster than interior counties: both the
binary 70-mile border indicator and the continuous log-distance measure are small and statistically
indistinguishable from no effect (hazard ratio $1.14$, $p=0.34$; hazard ratio $0.97$, $p=0.80$), consistent
with the static entry-density null -- sellers do not appear to race to be first into a border county
specifically. But conditional on a border county already having a first entrant, a \emph{second} dispensary
follows measurably faster the closer that county is to the nonlegal border: each additional log-mile of
distance lowers the hazard of a second entrant by roughly 30 percent (hazard ratio $0.70$, $p=0.003$),
though the corresponding binary measure is directionally consistent but not itself significant (hazard ratio
$1.26$, $p=0.16$). Figure~\ref{fig:entry_hazard_km} shows this pattern as Kaplan-Meier survival curves: the
border and interior curves for time to a first dispensary sit close together throughout, while the
second-dispensary curves visibly separate, with border counties' curve falling faster. We read this as
evidence that sellers do not treat a border county as an obviously better opportunity from the moment a
state's market opens -- consistent with the paper's other null entry-timing and entry-density results -- but
that once a border store reveals real demand, competitors respond to it faster than the same signal would
draw in an interior county. This is a narrower and more specific version of ``competing for the market''
than a simple race-to-enter story would predict, and it is direct behavioral evidence for a mechanism that
the revenue results alone can only imply.

\begin{figure}[htbp]
\centering
\includegraphics[width=0.95\textwidth]{figures/fig_entry_hazard_km.png}
\caption{Kaplan-Meier survival curves, border ($\le$70 miles) vs.\ interior counties, for time to a first
dispensary opening after a state's own retail launch (left) and time from a first to a second dispensary
(right), across the nine-state entry-density universe.}
\label{fig:entry_hazard_km}
\end{figure}

\subsubsection{Is the effect specific to the second entrant?}
The single rank-2 coefficient above raises an immediate question: is this something distinctive about a
market's \emph{second} entrant, or simply a generic ``more competition nearby $\to$ always faster entry''
pattern that would show up at every rank? We re-estimate the same hazard specification for the gap between
the second and third entrant, and between the third and fourth, among the subset of counties that reach
each milestone.

\input{tables/table_entry_hazard_by_rank}

Figure~\ref{fig:hazard_by_rank} plots these four hazard ratios together. The distance effect is significant
only at rank 2; it is essentially zero at rank 1, weakens and loses significance by rank 3, and is smaller
still by rank 4. This rules out a simple, monotonic ``agglomeration always speeds up entry'' story and is
consistent with the entry-threshold literature's general finding that competitive effects concentrate at low
entrant counts and are exhausted quickly \citep{BresnahanReiss1990,BresnahanReiss1991}, discussed in
Section~\ref{sec:lit_sequential}.

\begin{figure}[htbp]
\centering
\includegraphics[width=0.75\textwidth]{figures/fig_hazard_by_rank.png}
\caption{Hazard ratio on log(distance to nonlegal border), by entrant rank, with 95\% confidence intervals.
A ratio below 1 means faster entry (shorter time-to-event) closer to the border.}
\label{fig:hazard_by_rank}
\end{figure}

\subsubsection{Mechanism and placebo tests}
Two further checks are needed before treating the rank-2 result as evidence of a captive-demand mechanism
specifically. First, if the effect reflects the \emph{size} of the cross-border demand pool, it should
strengthen with the nonlegal-side population within reach, not just with raw proximity to the nearest
nonlegal county -- we construct a captive-population measure for each of the nine-state sample's 479
counties as the total population of \emph{all} nonlegal-state counties within 70 miles (not merely the
nearest one), via a radius query against national county population centroids; this is the same
construction intended for the paper's population-exposure measure (Section~4.7), applied here to the full
entry-timing sample rather than only Colorado and Washington. Second, if the effect is really about
\emph{nonlegal} borders and not a generic
feature of counties near any state line -- through-traffic, tourism, or other agglomeration unrelated to
cannabis policy -- it should not appear when we replace the nonlegal-border distance with distance to the
nearest county in \emph{any} other state, regardless of that state's legal status.

\input{tables/table_entry_hazard_mechanism}

The results split cleanly by test. The placebo test supports a nonlegal-border-specific interpretation: the
generic any-state distance measure is not significant on its own, loses to the nonlegal-border measure in a
horse race, and among counties whose nearest state-line neighbor happens to be a \emph{legal} state, distance
to that neighbor has essentially no relationship with second-entrant timing at all. The captive-population
test does not support a demand-size interpretation: the population-weighted measure is not significant
alone, is wrong-signed and insignificant when raced against distance, and is completely flat when the sample
is restricted to border counties only. Taken together, whatever is driving the rank-2 result responds to
\emph{proximity} to a nonlegal jurisdiction, not to the \emph{size} of the population behind it -- the
timing analogue of the travel-cost-threshold mechanism discussed in Section~\ref{sec:model_timing}, though
we cannot rule out that entrants simply use proximity as a coarse heuristic rather than a calibrated demand
estimate. The state-by-state breakdown in the same table shows that the pooled significance is concentrated
in Illinois, Massachusetts, and Ohio -- three states with the merit-lottery, host-community-agreement, and
staged-conversion licensing regimes described in Section~\ref{sec:institutions} -- while Colorado is
directionally consistent but not significant on its own, and Washington and Oregon show no effect at all.
This is a real limitation on how broadly the result can be attributed to this paper's core two-state revenue
sample specifically, and we return to it in Section~\ref{sec:discussion}.

\subsubsection{Does the equilibrium entry threshold move, or only the timing?}
A further, independent way to test the underlying idea is to ask whether cross-border demand shifts the
\emph{population} a county needs to support a second dispensary, using the classic entry-threshold approach
of \citet{BresnahanReiss1990} and \citet{BresnahanReiss1991} rather than a timing model. If a nearby captive
market genuinely supplements local demand, border counties should reach both the first- and second-store
population thresholds at a smaller local population than interior counties.

\input{tables/table_entry_threshold}

This prediction holds for the first store but not the second. Border counties do support a lone dispensary
at a smaller population than interior counties (35,436 versus 44,602), consistent with a demand supplement.
But the population increase needed to go from one store to two is \emph{larger}, not smaller, in border
counties (a threshold ratio of 1.67 versus 1.27) -- the opposite of what the demand-supplement story predicts
for the second-entrant margin specifically, even though that same margin is exactly where the timing effect
concentrates. We do not have a fully satisfying resolution of this tension and report it as an open question
rather than suppress it: entry \emph{speed} and the equilibrium entry \emph{threshold} appear to be
empirically dissociated in this setting, which is a narrower and more cautious conclusion than a clean
demand-threshold story would offer.

\subsection{Dose-Response in Nonlegal-Side Population Exposure}
Table~[INSERT TABLE NUMBER] reports the dose-response specification. The interaction between border
proximity (or continuous distance) and nonlegal-side population exposure is [INSERT SIGN, MAGNITUDE, AND
SIGNIFICANCE], indicating that the revenue premium in border counties rises as the population living on the
nonlegal side of the border grows. Figure~[INSERT FIGURE NUMBER] plots this relationship directly, showing
[INSERT DESCRIPTION, E.G. A BINNED SCATTER OF REVENUE PER DISPENSARY AGAINST NONLEGAL-SIDE POPULATION
EXPOSURE AMONG BORDER COUNTIES, RESIDUALIZED ON THE SAME CONTROLS AS THE REGRESSION]. This pattern indicates
that proximity to a nonlegal border alone does not fully explain the estimated revenue premium: dispensaries
that are both close to the border and accessible to a larger nonlegal-side population earn more per store
than dispensaries that are close to the border but face a small nonlegal-side population. We interpret this
result as mechanism evidence consistent with a cross-border demand channel, since a larger nearby
population represents a larger potential cross-border consumer market conditional on distance, rather than
as an independent causal confirmation of that channel.

\subsection{Placebo and Falsification Tests}
Table~[INSERT TABLE NUMBER] reports the placebo-sector results. Applying the same specification to
non-cannabis retail sectors -- specifically [INSERT SECTORS] -- does not produce a comparable border-
proximity pattern ($\hat\beta_1 = $ [INSERT ESTIMATE], [INSERT STANDARD ERROR] for [INSERT SECTOR]; [REPEAT
FOR EACH SECTOR]). This absence of a border effect in unrelated retail sectors is consistent with the
cannabis result reflecting demand specific to marijuana-policy differences across the border, rather than a
general tendency for all businesses near state borders to earn more revenue. As discussed in Section~2.4
and Section~4.8, we report the alcohol-sector placebo result separately and interpret it cautiously
($\hat\beta_1 = $ [INSERT ESTIMATE], [INSERT STANDARD ERROR]), because marijuana policy could plausibly
affect alcohol demand directly through substitution or complementarity \citep{BaggioChongKwon2020}; the
sectors with no plausible direct link to cannabis policy provide the cleaner falsification evidence.

\subsection{Summary of Main Findings}
Taken together, the results in this section indicate: (i) revenue per active dispensary is higher in legal-
state counties closer to a nonlegal-state border; (ii) this relationship is stronger where the nearby
nonlegal-side population is larger, producing a dose-response pattern; (iii) non-cannabis placebo sectors do
not display a comparable border pattern, with the partial exception of alcohol, which we interpret with
additional caution. [INSERT ANY ADDITIONAL SUMMARY POINTS ONCE FULL RESULTS ARE AVAILABLE.]

\section{Robustness}
\label{sec:robustness}

\subsection{Alternative Distance Thresholds and Continuous Distance}
We re-estimate the main specification using alternative border thresholds of [INSERT ALTERNATIVE
THRESHOLDS, E.G. A NARROWER AND A WIDER MILEAGE CUTOFF THAN THE MAIN SPECIFICATION] in place of the primary
threshold, and using the fully continuous log-distance specification described in Section~5.2. [INSERT
WHETHER THE RESULT IS STABLE ACROSS THRESHOLDS, OR WHETHER IT WEAKENS OR STRENGTHENS AT WIDER OR NARROWER
CUTOFFS, AND REPORT THE CORRESPONDING ESTIMATES.]

\subsection{Travel-Time Measures}
Where road-network or driving-time data are available, we replace straight-line distance with [INSERT
MEASURE, E.G. DRIVING TIME OR ROAD-NETWORK DISTANCE FROM A ROUTING SERVICE] and re-estimate the main and
dose-response specifications. [INSERT WHETHER RESULTS ARE SIMILAR TO THE STRAIGHT-LINE DISTANCE
SPECIFICATION, AND DISCUSS ANY MEANINGFUL DIFFERENCES, E.G. FOR COUNTIES WHERE ROAD ACCESS TO THE BORDER IS
INDIRECT.]

\subsection{Alternative Population-Exposure Buffers}
We recompute the nonlegal-side population-exposure measure using alternative radii of [INSERT ALTERNATIVE
RADII] in place of the main specification's radius, to test whether the dose-response result depends on a
particular choice of buffer distance. [INSERT WHETHER THE INTERACTION COEFFICIENT IS STABLE ACROSS
ALTERNATIVE BUFFERS.]

\subsection{Excluding Major Metropolitan and Tourism Counties}
Because large metropolitan counties and well-known tourist destinations could drive the main result for
reasons unrelated to cross-border cannabis demand, we re-estimate the main specification excluding [INSERT
CRITERION FOR EXCLUSION, E.G. THE LARGEST METROPOLITAN COUNTY IN EACH STATE, OR COUNTIES ABOVE A GIVEN
POPULATION OR TOURISM-INTENSITY THRESHOLD]. [INSERT WHETHER THE BORDER-PROXIMITY COEFFICIENT IS ROBUST TO
THIS EXCLUSION, AND REPORT THE CORRESPONDING ESTIMATE.]

\subsection{Controlling for Dispensary Entry and Licensing}
We add controls for [INSERT CONTROLS, E.G. THE NUMBER OF ACTIVE DISPENSARIES, AN INDICATOR OR CONTINUOUS
MEASURE OF STATE OR LOCAL LICENSING RESTRICTIVENESS, LOCAL MORATORIA OR BANS ON DISPENSARIES, AND YEARS
SINCE THE COUNTY'S OWN STATE LEGALIZED RETAIL SALES] to the main specification, to assess whether the
border-proximity coefficient reflects demand alone or is partly attributable to differences in local supply
conditions. [INSERT WHETHER THE COEFFICIENT IS STABLE, ATTENUATED, OR STRENGTHENED ONCE THESE CONTROLS ARE
ADDED.]

\subsection{Total County Revenue versus Revenue per Dispensary}
We report the main specification separately for (i) total county-level cannabis revenue, (ii) revenue per
active dispensary, (iii) the number of active dispensaries, and (iv) dispensary entry and exit, to
distinguish which margin drives any observed border effect, consistent with the identification discussion
in Section~5.6. [INSERT RESULTS FOR EACH OUTCOME AND DISCUSS WHETHER THE PATTERN IS CONSISTENT WITH A
SUPPLY-CONSTRAINED MARKET, AS THE CONCEPTUAL FRAMEWORK IN SECTION~3 WOULD PREDICT IF REVENUE PER DISPENSARY
RESPONDS MORE THAN DISPENSARY COUNT.]

\subsection{Non-Cannabis Placebo Sectors}
As reported in Section~6.3, we re-estimate the main specification for non-cannabis retail sectors. [INSERT
FULL SET OF PLACEBO-SECTOR ESTIMATES AND DISCUSS CONSISTENCY WITH THE CANNABIS-SPECIFIC INTERPRETATION.]

\subsection{Placebo Borders and Placebo Policy Dates}
As an additional falsification check, we construct placebo borders -- for example, borders between two
states that are both legal, or borders between two states that are both nonlegal, where no cannabis-policy
discontinuity exists -- and re-estimate the main specification using distance to these placebo borders in
place of distance to an actual legal-nonlegal border. [INSERT WHETHER A COMPARABLE PROXIMITY EFFECT APPEARS
AT PLACEBO BORDERS.] Where an event-study design is feasible (Section~5.4), we also assign placebo event
dates unrelated to any actual policy change and test whether a spurious break appears around those dates.
[INSERT RESULTS.]

\section{Discussion}
\label{sec:discussion}

The pattern of results described above, once finalized, would be consistent with a setting in which state-
level cannabis policy differences generate a meaningful fiscal externality across state lines. If residents
of a nonlegal state travel to a neighboring legal state to purchase cannabis, tax revenue that could
otherwise have accrued to the nonlegal state -- had it chosen to legalize and tax recreational sales itself
-- instead accrues to the neighboring state, along with the associated retail activity and any local
employment it supports. This is analogous to the interstate tax leakage documented in the cross-border
cigarette literature \citep{Lovenheim2008}, but with the added feature that the nonlegal state in this
setting is not simply taxing a legal good differently; it is prohibiting a good that a neighboring state has
chosen to permit and tax. The scale of this leakage, once measured directly rather than inferred, has
implications for how nonlegal states weigh the fiscal costs of continued prohibition against whatever
policy goals prohibition is intended to serve, and for how legal states weigh the fiscal benefits of
capturing this demand against the local costs -- traffic, enforcement burden, and any local externalities of
increased dispensary activity -- of hosting the counties that generate it.

These findings also bear on questions of interstate policy coordination. If neighboring states' cannabis
policies function partly like competing tax jurisdictions, then a state's decision to legalize or not
legalize is not fully independent of its neighbors' decisions: a nonlegal state's continued prohibition
generates identifiable revenue for legal-state border counties, which could itself be relevant information
in that state's own policy debate. Similarly, a legal state's decision about how tightly to restrict
dispensary licensing in its border counties is not purely a local question if border-county dispensaries
serve a substantially larger effective market than local population alone would suggest; underprovision of
licenses relative to that larger effective market could generate exactly the kind of revenue-per-dispensary
premium documented here without necessarily reflecting a well-functioning competitive market, an issue
directly connected to the market-power concerns raised in \citet{HollenbeckUetake2021}.

Several limitations should be kept in view when interpreting these results. Most importantly, the revenue
data used here do not identify the state of residence of individual consumers, so this paper cannot
directly attribute any specific share of border-county revenue to nonlegal-state residents; the population-
exposure and placebo-sector results provide indirect, mechanism-consistent evidence for a cross-border
demand interpretation, but they are not a substitute for direct data on consumer origin, such as survey
data, license-plate studies, or payment-card data, none of which are used in this paper. [INSERT ANY
ADDITIONAL DATA LIMITATIONS SPECIFIC TO THE FINAL SAMPLE, E.G. STATE COVERAGE GAPS OR DISCLOSURE
SUPPRESSION.] The paper's identification also relies primarily on cross-sectional and within-state variation
in border proximity, since a county's border status does not change over time; where the data permit an
event-study extension around a specific policy-timing event, this provides a stronger basis for causal
interpretation, but the core border-proximity and dose-response results should be read as strong
associations consistent with a cross-border demand mechanism rather than as fully causal estimates outside
of that event-study context.

\subsection{What the entry-timing result does and does not establish}
\label{sec:discussion_timing}

The entry-timing result in Section~\ref{sec:entry_hazard} deserves its own, more cautious accounting,
separate from the revenue-per-store result. It is a conditional statistical association -- a hazard ratio
from a stratified Cox model -- not a structural estimate of firm-level profits, demand elasticities, or a
genuine strategic-interaction game in the sense of \citet{Berry1992} or \citet{Seim2006}: we do not observe
firm identities, whether a second entrant actually observed or responded to the first, or what would have
happened under a counterfactual entry order. It passes one real test that strengthens a causal-flavored
reading -- a placebo comparison against generic (non-cannabis-relevant) border proximity -- but fails
another, since it does not show the corresponding shift in the classic cross-sectional entry threshold that
the same underlying story would predict, and it is statistically concentrated in three states (Illinois,
Massachusetts, Ohio) outside this paper's core two-state revenue sample. We think the right way to describe
this finding is as a \emph{secondary, mechanism-oriented} result that broadens the paper's evidence beyond
revenue levels to a second, independent margin -- entry timing -- rather than as a central contribution or
as a settled account of why cannabis markets structure themselves the way they do. The general regularity it
illustrates, that a market's second or third entrant matters disproportionately more than later ones, is
already well established in the entry literature \citep{BresnahanReiss1990,BresnahanReiss1991}; what is new
here is applying a timing-based design to that regularity in a licensed, border-adjacent retail setting, not
the regularity itself.

Three further tests, not conducted here, would help distinguish the competing explanations for this result
more directly, and we rank them by our own judgment of importance, feasibility with available or plausibly
obtainable data, and value to a reader trying to assess the mechanism. \emph{First} (high importance, high
feasibility): re-estimate the hazard models with an explicit interaction between border proximity and each
state's licensing-regime type (statewide lottery, local-option, merit-score-plus-lottery, staged
conversion), rather than only reporting separate by-state coefficients, to test directly whether the effect
requires the kind of sequential administrative clearance described in Section~\ref{sec:institutions}.
\emph{Second} (high importance, low feasibility with current data): test whether second-entrant timing
responds to the first entrant's \emph{realized performance} specifically -- e.g., its sales volume in the
months before the second entrant appears -- rather than merely its existence, which would more directly test
the information-revelation channel in Section~\ref{sec:model_timing}; this requires a store-level revenue
panel with opening dates that we do not currently have; state revenue data are collected at the county level.
\emph{Third} (moderate importance, high feasibility): test whether counties whose first and second licenses
were awarded in the same regulatory round (e.g., the same Washington or Illinois lottery cohort) show a
systematically different, more compressed gap than counties whose second entrant arrived through a separate,
later application round, which would help separate genuine sequential response from simultaneous
administrative allocation.

\section{Conclusion}
\label{sec:conclusion}

This paper examines whether counties in legal-cannabis states that are closer to the border of a nonlegal
state generate more cannabis revenue per active dispensary than counties farther from such a border, using
a national sample spanning multiple state borders rather than a single legal-nonlegal transition. Building
on \citet{HansenMillerWeber2020} and \citet{HaoCowan2020}, who each establish that cannabis policy gaps
across state borders generate real cross-border behavior in a single setting, we extend this evidence by
showing that the dispensary revenue premium associated with border proximity varies systematically with the
size of the population living on the nonlegal side nearby, and that no comparable pattern appears in
non-cannabis retail sectors. [INSERT ONE OR TWO SENTENCES SUMMARIZING THE FINAL HEADLINE MAGNITUDES ONCE
RESULTS ARE AVAILABLE.] We interpret these findings as consistent with a setting in which a meaningful share
of legal-state border-county dispensary revenue reflects demand from nearby residents of a neighboring
nonlegal state, though confirming this directly would require data on consumer residence that this paper
does not have. The results have implications for how both legal and nonlegal states weigh the fiscal and
policy consequences of state-level cannabis policy gaps, and for how legal states set dispensary licensing
policy in counties that appear to serve a substantially larger effective market than their own resident
population would suggest.

As a secondary contribution, we also show that a county's \emph{second} dispensary arrives systematically
faster the closer that county is to a nonlegal border, even though its first dispensary does not, and that
this pattern fades by the third and fourth entrant. This connects the paper to the entry-threshold tradition
in industrial organization \citep{BresnahanReiss1990,BresnahanReiss1991}, which establishes that a market's
second and third entrants matter disproportionately more than later ones; our contribution is to show a
timing-based version of that same regularity in a licensed, border-adjacent retail setting, using a hazard
model rather than the field's usual cross-sectional threshold approach. We are careful not to overstate this
result: it does not extend to a corresponding shift in the classic population threshold needed to support a
second store, and its statistical significance is concentrated in three states outside this paper's core
revenue sample. We report it as a mechanism-oriented finding that complements the revenue results, not as a
new theoretical contribution to the entry literature or as a central claim of the paper.

\appendix
\section{Extended Literature Review: Sequential Entry and the Economics of the Second Entrant}
\label{app:literature}

Section~\ref{sec:lit_sequential} summarizes briefly, in the main text, the literature most relevant to the
entry-timing result in Section~\ref{sec:entry_hazard}. Table~\ref{tab:app_literature} reports this
literature in more detail, including how each paper models entry, whether it distinguishes first, second,
and later entrants, and how closely it maps onto our own result. We searched systematically across
industrial organization, game theory, and applied microeconomics -- prioritizing the \textit{RAND Journal of
Economics}, \textit{Journal of Political Economy}, \textit{Econometrica}, \textit{Review of Economic
Studies}, \textit{American Economic Review}, and \textit{Quantitative Economics} -- and separately across
the cannabis-dispensary literature specifically. We verified every citation below against the publishing
journal or a repository record before inclusion; we found no paper, in economics or the cannabis-policy
literature, that documents or models a distinctive effect specifically associated with a market's second
dispensary.

\begin{table}[htbp]
\centering
\footnotesize
\caption{Extended Literature Summary}
\label{tab:app_literature}
\begin{tabular}{p{2.6cm}p{3.6cm}p{2.6cm}p{2.3cm}p{3.5cm}}
\toprule
Paper & Question / mechanism & Setting & 1st/2nd/later distinction? & Relation to our result \\
\midrule
\citet{BresnahanReiss1990} & How does the market size needed to support $N$ firms scale with $N$? &
  Isolated small-town markets (autos, etc.) & Yes, central object & Closest conceptual match: 2nd entry is
  qualitatively different from 1st and from 4th+, but via cross-sectional market size, not timing \\
\citet{BresnahanReiss1991} & Same question, five retail/professional industries & Doctors, dentists,
  druggists, plumbers, tire dealers & Yes, explicitly & Establishes that ``almost all variation in
  competitive conduct occurs with the entry of the second or third firm''; direct precedent for our
  rank-specific pattern \\
\citet{BerryReiss2007} & Survey of the entry-threshold literature & --- & Discusses generally & Background;
  confirms this is a settled general result \\
\citet{Berry1992} & Entry games with heterogeneous firms & Airline city-pairs & No, firms are heterogeneous
  rather than ordered & General framing; origin of the discrete-game entry literature alongside
  \citeauthor{BresnahanReiss1990} \\
\citet{Mazzeo2002} & Entry with endogenous product-quality choice & Highway motels & Yes, asks how a 2nd
  entrant's type choice responds to the 1st & Relevant if our 2nd dispensary differentiates from the 1st on
  price or quality; untested here \\
\citet{Seim2006} & Entry under incomplete information about rivals' locations & Video rental stores & Not by
  rank & Background on returns to differentiation \\
\citet{ToivanenWaterson2005} & How does a rival's presence affect a chain's own entry decisions? & UK burger
  duopoly, panel & Implicitly, via rival's lagged presence & Closest \emph{dynamic/timing} match: rival
  presence is a strong predictor of the other chain's subsequent entry \\
\citet{IgamiYang2016} & Dynamic multi-store entry with unobserved market heterogeneity & Canadian hamburger
  chains, 1970--2005 & Yes, sequential store-count dynamics & Methodological cousin (dynamic sequential
  entry), but models a firm's own multi-store rollout, not a rival's second store \\
\citet{GilbertVives1986} & Incumbents' joint incentive to deter entry & Theory & No & Background; deterrence,
  not sequencing \\
\citet{PrescottVisscher1977} & Sequential location with foresight & Theory & Yes, by construction & Framing
  device for the timing model in Section~\ref{sec:model_timing} \\
\citet{HollenbeckUetake2021} & Effect of WA's licensing restrictions on market power & WA dispensary
  licensing & No & Market power under a fixed license count, not sequencing; already central to our revenue
  discussion \\
\citet{DongTyndall2024} & Effect of dispensary presence on crime & WA 2014 license lottery & No & Uses the
  same WA lottery as our institutional discussion; no significant average crime effect \\
\citet{ThomasTian2021} & Effect of dispensary presence on property values & WA license lottery & No & Home
  prices fall 3--4\% within 0.36 miles of a new dispensary \\
\bottomrule
\end{tabular}
\end{table}

\newpage
\bibliographystyle{apalike}
\bibliography{references}

\end{document}
